%% LyX 2.4.4 created this file.  For more info, see https://www.lyx.org/.
%% Do not edit unless you really know what you are doing.
\documentclass[english,t]{beamer}
\usepackage{mathptmx}
\usepackage[T1]{fontenc}
\usepackage[latin9]{inputenc}
\usepackage{amssymb}
\usepackage{cancel}
\usepackage{graphicx}

\makeatletter

%%%%%%%%%%%%%%%%%%%%%%%%%%%%%% LyX specific LaTeX commands.
%% Because html converters don't know tabularnewline
\providecommand{\tabularnewline}{\\}
%% A simple dot to overcome graphicx limitations
\newcommand{\lyxdot}{.}


%%%%%%%%%%%%%%%%%%%%%%%%%%%%%% Textclass specific LaTeX commands.
% this default might be overridden by plain title style
\newcommand\makebeamertitle{\frame{\maketitle}}%
% (ERT) argument for the TOC
\AtBeginDocument{%
  \let\origtableofcontents=\tableofcontents
  \def\tableofcontents{\@ifnextchar[{\origtableofcontents}{\gobbletableofcontents}}
  \def\gobbletableofcontents#1{\origtableofcontents}
}

%%%%%%%%%%%%%%%%%%%%%%%%%%%%%% User specified LaTeX commands.
\usetheme{Warsaw}
% or ...

\setbeamercovered{transparent}
% or whatever (possibly just delete it)
\setbeamercovered{invisible} %makes overlays invisible

%gets rid of navigation symbols
\setbeamertemplate{navigation symbols}{}

%\usepackage{hyperref}
%\usepackage[usenames,dvipsnames,table]{xcolor} %adds additional color options
\definecolor{links}{HTML}{2A1B81}
%\hypersetup{colorlinks,linkcolor=blue,urlcolor=links}

\usepackage{multicol}

\usepackage{tikz}
\usetikzlibrary{calc}
\usetikzlibrary{plotmarks}
\usetikzlibrary{shapes}
\usepackage{pgfplots}
\pgfplotsset{compat=newest}

\usepackage{forest} %%for tree diagram

\usepackage{calc}
\usepackage[nomessages]{fp}

%\usepackage{advdate}    % Advancing/saving dates
\usepackage{datetime}   % Dates formatting
%\usepackage{datenumber} % Counters for dates

\usepackage[super]{nth}

\usepackage{etex}

\usepackage{fancybox}

%%table colors
%\usepackage[table]{xcolor}
\usepackage{colortbl}

\usepackage[outline]{contour}  %allows text halos
\contourlength{1pt} %sets halo width
%%example: \node at (0.75,0.5) {\contour{white}{\Large Text!}};
%%example:  \node [fill=white,inner sep=1pt] at (2.25,0.5) {\Large 			Text!};
%%example:  \node [fill=white,rounded corners=2pt,inner sep=1pt] at 			(3.75,0.5) {\Large Text!};

%%%%%%%%%%%%%%%%%%%%%%%%%%%
\def\formatdateny{\csname noyear\languagename\expandafter\endcsname\formatdate}
\def\noyearenglish#1, \the\@year{#1}
\let\noyearamerican\noyearenglish
\let\noyearbritish\noyearenglish

%%allows uncover with forest diagrams
\tikzset{
    invisible/.style={opacity=0,text opacity=0},
    visible on/.style={alt=#1{}{invisible}},
    alt/.code args={<#1>#2#3}{%
      \alt<#1>{\pgfkeysalso{#2}}{\pgfkeysalso{#3}} % \pgfkeysalso doesn't change the path
    },
}
\forestset{
  visible on/.style={
    for tree={
      /tikz/visible on={#1},
      edge={/tikz/visible on={#1}}}}}

%%%redefines column alignment to match vertically
%\define@key{beamerframe}{t}[true]{% top
  %\beamer@frametopskip=.2cm plus .5\paperheight\relax%
  %\beamer@framebottomskip=0pt plus 1fill\relax%
 % \beamer@frametopskipautobreak=\beamer@frametopskip\relax%
  %\beamer@framebottomskipautobreak=\beamer@framebottomskip\relax%
  %\def\beamer@initfirstlineunskip{}%
%}

%%provides pdf with 4 slides per page; don't forget to add 'handout' to remove overlays
%\usepackage{pgfpages}
%\pgfpagesuselayout{4 on 1}[a4paper, landscape, border shrink=7mm]
%\pgfpageslogicalpageoptions{1}{border code=\pgfusepath{stroke}}
%\pgfpageslogicalpageoptions{2}{border code=\pgfusepath{stroke}}
%\pgfpageslogicalpageoptions{3}{border code=\pgfusepath{stroke}}
%\pgfpageslogicalpageoptions{4}{border code=\pgfusepath{stroke}}
%%%Don't forget to add 'handout'

\makeatother

\usepackage{babel}
\begin{document}
\newcommand{\xmax}{2000}
\newcommand{\xmin}{0}
\newcommand{\xstep}{0} %must be xmin+n
\newcommand{\ymax}{500}
\newcommand{\ymin}{0}
\newcommand{\ystep}{0} %must be ymin+n

\newcounter{countEx} \setcounter{countEx}{0}
\newcounter{countProb} \setcounter{countProb}{0}
\resetcounteronoverlays{countEx} \resetcounteronoverlays{countProb}

\newcommand{\ex}[3]{
	\addtocounter{countEx}{1}
	\only<#1-#2>{\begin{exampleblock} {Example \chNum.\secNum.\arabic{countEx}.} #3	\end{exampleblock}}
}

\newcommand{\prob}[4]{
	\addtocounter{countProb}{1}
	\only<#1-#2>{\begin{exampleblock} {Problem  \chNum.\secNum.\arabic{countProb}.\,#3} #4	\end{exampleblock}}
}

\newcommand{\yline}[4]{
	(#2 - #4)/(#1 - #3)*(x - #1) + #2
}

\beamerdefaultoverlayspecification{<+->}

%%%%Chapter and Section numbers%%%%%
\newcommand{\chNum}{3}
\newcommand{\secNum}{2}
\newcommand{\secTitle}{Basic Terms of Probability}
\date{\formatdate{3}{10}{2014}}

\title[Section \chNum.\secNum: \secTitle]{Section \chNum.\secNum: \secTitle}

\subtitle{Finite Mathematics~~MTH 107~~Fall 2014}

\author{Dr.\,James Ashe}

\titlegraphic{\includegraphics[height=2.75cm]{/home/jim/JamesAshe/MHU/images/MarsHilUniversityLogo-Virtical.png}}

\institute{Mars Hill University}

\makebeamertitle 

%%True false command
\newcommand{\T}{\textcolor{green}{\contour{black}{T}}}
\newcommand{\F}{\textcolor{red}{\contour{black}{F}}}
\newcommand{\uT}[2]{\uncover<#1-#2>{\T}}
\newcommand{\uF}[2]{\uncover<#1-#2>{\F}}

%tkiz ball item 
\newcommand*\circled[1]{\tikz[baseline=(char.base)]
	{\node[circle,ball color=green!50!black!100, shade,   color=white,inner sep=1.2pt] (char) {\tiny #1};
	}
}
\begin{frame}{Set Theory Notation}

\begin{block}{Recall:}

\begin{itemize}
\item []<2->We use capital letters to denote sets, and declare its elements
with brackets.

\begin{itemize}
\item [ex.]<3->$E=\{3,4,5,6,7\}$ is the set of integers between $3$ and
$7$.
\end{itemize}
\item []<4->We use $n(E)$ to denote the \emph{number of elements in $E$.}

\begin{itemize}
\item [ex.]<5->$n(E)=5$
\end{itemize}
\item []<6->We use an accent to denote that a set is the \emph{complement}.

\begin{itemize}
\item [ex.]<7->Let $S=\{1,2,3,4,5,6,7\}$ be the universal set of $E$.
Then

\begin{itemize}
\item []<8->$E'=\{1,2\}$
\end{itemize}
\end{itemize}
\end{itemize}
\end{block}
\end{frame}

\begin{frame}{Basic Terms}

\vspace{-.5cm}
\begin{columns}


\column{5.5cm}
\begin{exampleblock}{}
A die is rolled. Find:

\begin{enumerate}
\item [a.]<2->The sample space
\item [b.]<4->The event ``rolling a 4''
\item [c.]<7->The probability of rolling a 4
\item [d.]<14->The odds of ``rolling a 4''
\item [e.]<21->The probability of rolling below 4
\item [f.]<23->The odds of rolling a below 4
\end{enumerate}
\end{exampleblock}

\column{5.5cm}

\vspace{5pt}\emph{Solution}
\begin{enumerate}
\item [a.]<3->$S=\{1,2,3,4,5,6\}$
\item [b.]<6->$E_{1}=\{4\}$
\item [c.]<12->$p(E_{1})$\uncover<13->{\textrm{$=\frac{n(E_{1})}{n(S)}=\frac{1}{6}$}}
\item [d.]<19->$o(E_{1})$\uncover<20->{\textrm{$=n(E_{1}):n(E_{1}')=1:5$}}
\item [e.]<22->$E_{2}=\{1,2,3\}$. $p(E_{2})=\frac{n(E_{2})}{n(S)}=\frac{3}{6}=\frac{1}{2}$
\item [f.]<24->$o(E_{2})=3:3=1:1$
\end{enumerate}
\end{columns}
\only<1-7>{
\begin{block}{Basic Probability Terms}

\uncover<1->{\textbf{Experiment: }a process by which an outcome is
obtained.}

\uncover<2->{\textbf{Sample space: }set $S$ of all possible outcomes
of an experiment.}

\uncover<5->{\textbf{Event: }any subset $E\subseteq S$.}
\end{block}
}

\only<8-14>{
\begin{block}{Probability of an event, \$E\$}
The \emph{probability} of event $E$, denoted $p(E)$, is

\hfill \uncover<9->{$p(E)=\begin{array}{c}
\underline{{n(E)}}\\
n(S)
\end{array}$}%
\begin{tabular}{l}
\uncover<10->{\textcolor{red}{\footnotesize $\leftarrow$ number of ways $E$ can happen}}\tabularnewline
\uncover<11->{\textcolor{red}{\footnotesize $\leftarrow$ total number of outcomes}}\tabularnewline
\end{tabular}\hfill

if the experiment's outcomes are equally likely.
\end{block}
}\only<15->{
\begin{block}{Odds of an event, \$E\$}
The \emph{odds} of event $E$, denoted $o(E)$, are

\uncover<16->{%
\begin{tabular}{crc}
 & $o(E)=n(E):n(E')$\uncover<18->{\textcolor{red}{\footnotesize $\leftarrow$ number of ways $E$ CANNOT happen}} & \tabularnewline
 & \uncover<17->{\textcolor{red}{\footnotesize  $\uparrow$ number of ways $E$ can happen}} \hspace{2.5cm} & \tabularnewline
\end{tabular}}

if the experiment's outcomes are equally likely.
\end{block}
}
\end{frame}

\begin{frame}{Relative Frequency}

\begin{columns}


\column{5.5cm}

\textbf{Probability }

\ex{2}{}{If a students is randomly selected from the class, what
is the probability that a female is selected?}
\begin{itemize}
\item <3->Theoretical
\end{itemize}

\column{5.5cm}

\textbf{Relative Frequency}

\ex{4}{}{A baseball player has $2,865$ hits in $8,239$ times at
bat. What are the chances of him getting a hit at his next attempt?}
\begin{itemize}
\item <5->Empirical
\end{itemize}
\end{columns}
\begin{block}<6->{Law of Large Numbers}
If an experiment is repeated a large number of times, the relative
frequency of an outcome will become close the actual probability of
that outcome. 
\end{block}
\end{frame}

\begin{frame}{Relative Frequency}

\begin{figure}
\includegraphics[width=7cm]{../../../../.wuala/Data/Temp/Export/pics/400px-Largenumbers.svg.png}

\caption{As the number of rolls increase, the average value approaches $3.5$}

\end{figure}

\end{frame}

\begin{frame}{Problems}

\prob{1}{1}{A jar contains $12$ black, $8$ red, $10$ yellow, and $5$ green jelly beans. Find the probability of the following:}{The probability that you randomly pick a black jelly bean.}\prob{2}{2}{A jar contains $12$ black, $8$ red, $10$ yellow, and $5$ green jelly beans. Find the following:}{The odds that you randomly pick a black jelly bean.}\prob{3}{3}{A jar contains $12$ black, $8$ red, $10$ yellow, and $5$ green jelly beans. Find the following:}{The probability that you randomly pick a black or red jelly bean.}\prob{4}{4}{A jar contains $12$ black, $8$ red, $10$ yellow, and $5$ green jelly beans. Find the following:}{The odds that you randomly pick a black or red jelly bean.}\prob{5}{5}{A jar contains $12$ black, $8$ red, $10$ yellow, and $5$ green jelly beans.Find the following:}{The probability that you randomly pick a black and red jelly bean.}\prob{6}{6}{A jar contains $12$ black, $8$ red, $10$ yellow, and $5$ green jelly beans. Find the following:}{The odds that you randomly pick a jelly bean that is not green.}

\prob{7}{7}{A jar contains $12$ black, $8$ red, $10$ yellow, and $5$ green jelly beans. Find the following:}{You randomly pick a jelly bean. What is the sample space?}

\prob{8}{8}{Find the following:}{What is the probability of flipping a coin 3 times and getting all heads?}\prob{9}{9}{Find the following:}{What is the probability of flipping a coin 3 times and getting no heads?}\prob{10}{10}{Find the following:}{What is the probability of flipping a coin 3 times and getting two heads?}\prob{11}{11}{Find the following:}{What is the probability of flipping a coin 3 times and getting at least two heads?}

\prob{12}{12}{Find the following:}{What is the probability of flipping a coin 3 times and getting one tail?}

\prob{13}{13}{Find the following:}{What is the probability of flipping a coin 1 time and getting a tail \emph{or} a head?}

\prob{14}{14}{Find the following:}{What is the probability of flipping a coin 1 time and getting a tail \emph{and} a head?}

\prob{15}{15}{Find the following:}{A card is drawn from a 52 card deck. What is the probability of getting a queen?}

\prob{16}{16}{Find the following:}{A card is drawn from a 52 card deck. What is the probability of getting a queen of hearts?}

\prob{17}{17}{Find the following:}{A card is drawn from a 52 card deck. What is the probability of getting a queen or a heart?}
\end{frame}

\end{document}
